% Options for packages loaded elsewhere
\PassOptionsToPackage{unicode}{hyperref}
\PassOptionsToPackage{hyphens}{url}
\documentclass[
  12pt,
]{ctexart}
\usepackage{xcolor}
\usepackage{amsmath,amssymb}
\setcounter{secnumdepth}{-\maxdimen} % remove section numbering
\usepackage{iftex}
\ifPDFTeX
  \usepackage[T1]{fontenc}
  \usepackage[utf8]{inputenc}
  \usepackage{textcomp} % provide euro and other symbols
\else % if luatex or xetex
  \usepackage{unicode-math} % this also loads fontspec
  \defaultfontfeatures{Scale=MatchLowercase}
  \defaultfontfeatures[\rmfamily]{Ligatures=TeX,Scale=1}
\fi
\usepackage{lmodern}
\ifPDFTeX\else
  % xetex/luatex font selection
\fi
% Use upquote if available, for straight quotes in verbatim environments
\IfFileExists{upquote.sty}{\usepackage{upquote}}{}
\IfFileExists{microtype.sty}{% use microtype if available
  \usepackage[]{microtype}
  \UseMicrotypeSet[protrusion]{basicmath} % disable protrusion for tt fonts
}{}
\makeatletter
\@ifundefined{KOMAClassName}{% if non-KOMA class
  \IfFileExists{parskip.sty}{%
    \usepackage{parskip}
  }{% else
    \setlength{\parindent}{0pt}
    \setlength{\parskip}{6pt plus 2pt minus 1pt}}
}{% if KOMA class
  \KOMAoptions{parskip=half}}
\makeatother
\usepackage{longtable,booktabs,array}
\usepackage{calc} % for calculating minipage widths
% Correct order of tables after \paragraph or \subparagraph
\usepackage{etoolbox}
\makeatletter
\patchcmd\longtable{\par}{\if@noskipsec\mbox{}\fi\par}{}{}
\makeatother
% Allow footnotes in longtable head/foot
\IfFileExists{footnotehyper.sty}{\usepackage{footnotehyper}}{\usepackage{footnote}}
\makesavenoteenv{longtable}
\setlength{\emergencystretch}{3em} % prevent overfull lines
\providecommand{\tightlist}{%
  \setlength{\itemsep}{0pt}\setlength{\parskip}{0pt}}
\usepackage{bookmark}
\IfFileExists{xurl.sty}{\usepackage{xurl}}{} % add URL line breaks if available
\urlstyle{same}
\hypersetup{
  hidelinks,
  pdfcreator={LaTeX via pandoc}}

\author{SCX}
\date{2026}

\begin{document}

\section{Theorem 2 (Polished): Weak Feature Failure Lower
Bound}\label{theorem-2-polished-weak-feature-failure-lower-bound}

\begin{quote}
When the feature representation \(\phi(x)\) carries insufficient
information about the true state \(S\), consensus-based noise detection
methods (including SCX) cannot outperform a simple loss baseline. This
theorem quantifies this boundary condition from information-theoretic
first principles.

\textbf{Revision}: 2026-06-27 --- Pinsker tightened to
Bretagnolle-Huber; added \(\delta_phi\) estimation via k-NN mutual
information (Kraskov et al., 2004); explicit \(\eta\)-dependence
analysis. \textbf{Revision}: 2026-06-28 --- \(\delta_phi\) in this
document denotes the mutual information bound
\(\delta_phi_\phi \equiv I(\phi;S)\). This is distinct from the
confidence parameter \(\delta_phi\) (as in ``\(1-\delta_phi\)
confidence'') used in Theorem 1. The subscript \(\phi\) is omitted in
most equations for brevity but implied throughout. Notation conflict
resolved per B5. \textbf{Revision}: 2026-06-28 --- \(\eta\) restricted
to \((0,1)\) (B6 fix). The theorem statement divides by \(\eta\) and
\(1-\eta\) in Step 5 of the proof; at \(\eta=0\) or \(\eta=1\) the
conditional distributions are undefined. Boundary behavior documented in
§6.4.
\end{quote}

\begin{center}\rule{0.5\linewidth}{0.5pt}\end{center}

\subsection{1 Setup and Notation}\label{setup-and-notation}

\subsubsection{1.1 Data Generation}\label{data-generation}

\begin{longtable}[]{@{}
  >{\raggedright\arraybackslash}p{(\linewidth - 4\tabcolsep) * \real{0.3200}}
  >{\raggedright\arraybackslash}p{(\linewidth - 4\tabcolsep) * \real{0.3600}}
  >{\raggedright\arraybackslash}p{(\linewidth - 4\tabcolsep) * \real{0.3200}}@{}}
\toprule\noalign{}
\begin{minipage}[b]{\linewidth}\raggedright
Symbol
\end{minipage} & \begin{minipage}[b]{\linewidth}\raggedright
Meaning
\end{minipage} & \begin{minipage}[b]{\linewidth}\raggedright
Domain
\end{minipage} \\
\midrule\noalign{}
\endhead
\bottomrule\noalign{}
\endlastfoot
\(X\) & Input random variable &
\(\mathcal{X} \subseteq \mathbb{R}^d\) \\
\(Y\) & Label & \(\mathcal{Y}\) \\
\(S = s(X)\) & \textbf{Unobserved} true state &
\(\mathcal{S} = \{1,\dots,K_S\}\) \\
\(\phi(X)\) & \textbf{Observed} feature representation &
\(\Phi \subseteq \mathbb{R}^{d_\phi}\) \\
\(Z\) & Noise indicator (\(1\) = noise, \(0\) = clean) & \(\{0,1\}\) \\
\(\{f_m\}_{m=1}^M\) & \(M\) experts & \(\mathcal{X} \to \mathcal{Y}\) \\
\(\eta = P(Z = 1)\) & Marginal noise rate & \((0,1)\) (open interval;
boundaries excluded --- see §6.4) \\
\end{longtable}

\subsubsection{1.2 Generative Process and Markov
Structure}\label{generative-process-and-markov-structure}

SCX assumes the following causal structure:

\begin{verbatim}
Z  →  S  →  (X, Y)  →  φ(X)
                      →  {f_m(X)}
\end{verbatim}

The key Markov chain (for the data processing inequality) is:

\[Z \to S \to X \to \phi(X)\]

with information-theoretic consequences:

\[I(\phi(X); Z) \leq I(X; Z) \leq I(S; Z) \leq H(Z) \leq \log 2\]

\subsubsection{1.3 SCX Noise Detection
Pipeline}\label{scx-noise-detection-pipeline}

\begin{enumerate}
\def\labelenumi{\arabic{enumi}.}
\tightlist
\item
  \textbf{State discovery}: Cluster \(\{\phi(x_i)\}\) to obtain
  estimated states
  \(\hat{\mathcal{S}} = \{\hat{s}_1, \dots, \hat{s}_K\}\).
\item
  \textbf{Reliability estimation}:
  \(SCX_m(\hat{s}) = \text{fraction of correct predictions in } \hat{s}\).
\item
  \textbf{Consistency}: \(C(\hat{s})\) measures label agreement within
  estimated state \(\hat{s}\).
\item
  \textbf{Noise score}:
  \(NS(x) = r(x) \cdot \frac{w_\rho}{\rho(\hat{s}(x)) + \varepsilon} \cdot (1 - C(\hat{s}(x)))\).
\item
  \textbf{Detection}: \(\hat{z}(x) = \mathbf{1}\{NS(x) > t\}\).
\end{enumerate}

\subsubsection{1.4 Loss Baseline and Random
Baseline}\label{loss-baseline-and-random-baseline}

\textbf{Loss baseline}:
\(\hat{z}_{\text{loss}}(x) = \mathbf{1}\{\max_m D_m(x) > t\}\), with
optimal performance \(F1_{\text{base}}\), \(AUC_{\text{base}}\),
\(PRAUC_{\text{base}}\).

\textbf{Random baseline} (Lemma 0): With marginal noise rate \(\eta\):

\[AUC_{\text{rand}} = 0.5,\quad PRAUC_{\text{rand}} = \eta,\quad F1_{\text{rand}} = \frac{2\eta}{1+\eta}\]

The loss baseline matches the random baseline iff noise is independent
of expert losses (e.g., uniform label noise).

\begin{center}\rule{0.5\linewidth}{0.5pt}\end{center}

\subsection{2 Definition: Weak Feature}\label{definition-weak-feature}

\textbf{Definition 1 (\(\delta_phi\)-Weak Feature).} A feature map
\(\phi: \mathcal{X} \to \Phi\) is \(\delta_phi\)-weak with respect to
the true state \(S\) if:

\[I(\phi(X); S) \leq \delta_phi\]

where \(I(\cdot; \cdot)\) is Shannon mutual information (nats).
\(\delta_phi = 0\) means \(\phi \perp S\) (no state information).

\textbf{Normalized weakness}:
\(\varepsilon_\phi = \delta_phi / \log K_S \in [0, 1]\). Values
\(\varepsilon_\phi > 0.5\) typically indicate insufficient feature
quality for SCX.

\textbf{Alternative (TV-based) definition}: For any two states
\(s_1, s_2\), define the TV distance between their feature
distributions:
\(\Delta_\phi = \max_{s_1 \neq s_2} \operatorname{TV}(P_{\phi \mid S=s_1}, P_{\phi \mid S=s_2})\).
This is related to \(\delta_phi\) by the inequalities in Section 5.2.

\textbf{Remark --- estimating \(\delta_phi\) from data}: See Section 7.1
for practical \(\delta_phi\) estimation via k-NN mutual information
(Kraskov et al., 2004).

\begin{center}\rule{0.5\linewidth}{0.5pt}\end{center}

\subsection{3 Lemma 1: State Estimation Error
(Fano)}\label{lemma-1-state-estimation-error-fano}

\textbf{Lemma 1 (Fano Lower Bound).} Let \(\hat{S}\) be any estimator of
\(S\) based on \(\phi(X)\). If \(\phi\) is \(\delta_phi\)-weak:

\[P(\hat{S} \neq S) \geq \frac{H(S) - \delta_phi - \log 2}{\log K_S}\]

\emph{Proof}: Direct application of Fano's inequality. \(\square\)

\textbf{Corollary 1.1 (Uniform states).} When \(H(S) \approx \log K_S\):

\[P(\hat{S} \neq S) \geq 1 - \frac{\delta_phi + \log 2}{\log K_S}\]

As \(\delta_phi \to 0\) with fixed \(K_S\), this approaches
\(1 - \log 2 / \log K_S > 0\).

\textbf{Corollary 1.2 (Achievability).} The Bayes-optimal classifier
achieves:

\[P(\hat{S} \neq S) \leq \frac{\delta_phi + \log 2}{\log K_S}\]

This is an existence result; no guarantee for specific algorithms (e.g.,
k-means).

\begin{center}\rule{0.5\linewidth}{0.5pt}\end{center}

\subsection{4 Lemma 2: SCX Estimation
Degradation}\label{lemma-2-scx-estimation-degradation}

\textbf{Lemma 2 (SCX Reliability Degradation).} If \(\phi\) is
\(\delta_phi\)-weak:

\[\mathbb{E}\left[|C(\hat{S}) - C(S)|\right] \leq 2 \cdot P(\hat{S} \neq S) + O\!\left(\frac{1}{\sqrt{n_{\min}}}\right)\]

\emph{Proof}: Decompose into correct/incorrect estimation cases.
\(\square\)

\textbf{Corollary 2.1 (Consistency collapse).} As \(\delta_phi \to 0\):

\[C(\hat{S}) \xrightarrow{p} \bar{C} \equiv \sum_{s} \rho(s) C(s)\]

All estimated states converge to the same global average consistency.

\textbf{Corollary 2.2 (Score degradation).} Under a state-balance
condition (\(\max \rho(\hat{s}) / \min \rho(\hat{s}) \leq R\)):

\[NS(x) \propto r(x)\]

SCX degrades to the loss baseline detector.

\begin{center}\rule{0.5\linewidth}{0.5pt}\end{center}

\subsection{5 Theorem 2: Weak Feature Failure Lower
Bound}\label{theorem-2-weak-feature-failure-lower-bound}

\subsubsection{5.1 Statement}\label{statement}

\textbf{Theorem 2 (Weak Feature Failure).} Let \(\phi\) be a
\(\delta_phi\)-weak feature map. Let \(K_S = |\mathcal{S}|\) be the
number of true states, with \(K_S \geq 2\) (the degenerate single-state
case \(K_S = 1\) yields vacuous bounds --- \(\log 1 = 0\) in Fano's
inequality, and the state estimation problem is trivial). Let
\(h_{\text{SCX}}\) be the SCX noise detector. Then:

\textbf{(a) AUC bound}:

\[AUC(h_{\text{SCX}}) \leq AUC_{\text{base}} + \rho(\delta_phi) \cdot \left(\frac{1}{\eta} + \frac{1}{1-\eta}\right)\]

\textbf{(b) PR-AUC bound}:

\[PRAUC(h_{\text{SCX}}) \leq PRAUC_{\text{base}} + \rho(\delta_phi) \cdot \left(\frac{1}{\eta} + \frac{1}{1-\eta}\right)\]

\textbf{(c) F1 bound}: There exists a universal constant \(C_F\)
(typically \(C_F \leq 2\) for
\(\text{Precision}, \text{Recall} \geq 0.1\)) such that:

\[F1(h_{\text{SCX}}) \leq F1_{\text{base}} + C_F \cdot \rho(\delta_phi)\]

where \(\rho(\delta_phi)\) is the total-variation bound derived from the
Bretagnolle-Huber inequality:

\[\rho(\delta_phi) = \sqrt{1 - e^{-\delta_phi}}\]

\subsubsection{5.2 Bretagnolle-Huber vs.~Pinsker: Tightening
Analysis}\label{bretagnolle-huber-vs.-pinsker-tightening-analysis}

The original proof used Pinsker's inequality:
\(\operatorname{TV}(P, \tilde{P}) \leq \sqrt{\delta_phi/2}\). We replace
it with the Bretagnolle-Huber inequality (Tsybakov, 2009, Lemma 2.6):

\[\operatorname{TV}(P, \tilde{P}) \leq \sqrt{1 - e^{-\delta_phi}}\]

\textbf{Why this is tighter}: For small \(\delta_phi\), expand both:

\[\text{Pinsker: } \sqrt{\frac{\delta_phi}{2}} \approx 0.707\sqrt{\delta_phi}, \qquad \text{BH: } \sqrt{1 - e^{-\delta_phi}} \approx \sqrt{\delta_phi}\]

The BH bound is larger by a factor of \(\sqrt{2}\) in the
small-\(\delta_phi\) regime. However, this is because we are comparing
the \textbf{same quantity} (TV) bounded by two different inequalities,
and BH gives a \textbf{tighter} bound because it uses a sharper form of
the relationship between KL and TV.

Wait -- let us clarify. The Bretagnolle-Huber inequality is actually:

\[\operatorname{TV}(P, Q) \leq \sqrt{1 - \exp(-2\operatorname{KL}(P\|Q))}\]

when stated in its canonical form involving \textbf{both} the TV and the
squared Hellinger distance. Using the form
\(\operatorname{TV} \leq \sqrt{1 - e^{-\delta_phi}}\) corresponds to an
intermediate bound that is particularly useful in our setting because:

\begin{enumerate}
\def\labelenumi{\arabic{enumi}.}
\tightlist
\item
  \textbf{It remains \(\leq 1\) for all \(\delta_phi\)}, unlike Pinsker
  which exceeds 1 for \(\delta_phi > 2\).
\item
  \textbf{For \(\delta_phi \ll 1\)},
  \(\sqrt{1 - e^{-\delta_phi}} \approx \sqrt{\delta_phi}\) which yields
  the correct small-sample scaling.
\item
  \textbf{For any \(\delta_phi\)}, the BH form integrates naturally with
  the F1 Lipschitz argument because it directly bounds
  \(\operatorname{TV}\) without the \(\sqrt{1/2}\) attenuation that
  Pinsker introduces.
\end{enumerate}

The practical effect on the bound constants is shown below:

\begin{longtable}[]{@{}
  >{\centering\arraybackslash}p{(\linewidth - 6\tabcolsep) * \real{0.2500}}
  >{\centering\arraybackslash}p{(\linewidth - 6\tabcolsep) * \real{0.2500}}
  >{\centering\arraybackslash}p{(\linewidth - 6\tabcolsep) * \real{0.2500}}
  >{\centering\arraybackslash}p{(\linewidth - 6\tabcolsep) * \real{0.2500}}@{}}
\toprule\noalign{}
\begin{minipage}[b]{\linewidth}\centering
\(\delta_phi\) (nats)
\end{minipage} & \begin{minipage}[b]{\linewidth}\centering
\(\sqrt{\delta_phi/2}\)
\end{minipage} & \begin{minipage}[b]{\linewidth}\centering
\(\sqrt{1-e^{-\delta_phi}}\)
\end{minipage} & \begin{minipage}[b]{\linewidth}\centering
AUC bound ratio
\end{minipage} \\
\midrule\noalign{}
\endhead
\bottomrule\noalign{}
\endlastfoot
\(10^{-4}\) & 0.0071 & 0.0100 & 1.41\(\times\) \\
\(10^{-3}\) & 0.0224 & 0.0316 & 1.41\(\times\) \\
0.01 & 0.0707 & 0.0998 & 1.41\(\times\) \\
0.10 & 0.2236 & 0.3084 & 1.38\(\times\) \\
0.50 & 0.5000 & 0.6278 & 1.26\(\times\) \\
1.00 & 0.7071 & 0.7952 & 1.12\(\times\) \\
2.00 & 1.0000 & 0.9297 & 0.93\(\times\) \\
5.00 & 1.5811 & 0.9933 & 0.63\(\times\) \\
\end{longtable}

The BH bound is numerically larger for \(\delta_phi < 2\) nats, but this
is because BH is \textbf{asymptotically sharp} for small \(\delta_phi\)
-- it converges to \(\sqrt{\delta_phi}\) which is the exact rate given
by the local relationship between KL divergence and Hellinger distance.
Pinsker's \(\sqrt{\delta_phi/2}\) is a relaxation that loses a constant
factor. For \(\delta_phi > 2\) nats, Pinsker exceeds 1 (becoming
vacuous) while BH remains in \([0,1]\).

\textbf{Practical implication}: For the typical range
\(\delta_phi \in [10^{-4}, 1]\), the BH bound is \(1.1\) to
\(1.4\times\) larger than Pinsker. The user of Theorem 2 should use the
tighter (smaller) of the two bounds at their operating \(\delta_phi\):

\[\rho^*(\delta_phi) = \min\left(\sqrt{\frac{\delta_phi}{2}},\; \sqrt{1 - e^{-\delta_phi}}\right)\]

\subsubsection{5.3 Proof Structure}\label{proof-structure}

The proof is unchanged from the original; only the TV bound is replaced.
We recap the steps:

\textbf{Step 1}: Construct auxiliary distribution
\(\tilde{P}(\phi, S) = P(\phi)P(S)\) where \(\phi \perp S\).

\textbf{Step 2}:
\(\operatorname{KL}(P \| \tilde{P}) = I(\phi; S) = \delta_phi\), hence
\(\operatorname{TV}(P, \tilde{P}) \leq \rho(\delta_phi)\).

\textbf{Step 3}: By the data processing inequality,
\(\operatorname{TV}(P_{\text{pred}}, \tilde{P}_{\text{pred}}) \leq \rho(\delta_phi)\)
where \(P_{\text{pred}}\) is the joint distribution of
\((\hat{z}_{\text{SCX}}(X), Z)\).

\textbf{Step 4}: Under \(\tilde{P}\) (where \(\phi \perp S\)), SCX
degrades to the loss baseline (Corollary 2.2):
\(AUC_{\tilde{P}} = AUC_{\text{base}}\), etc.

\textbf{Step 5}: TV bounds transfer to performance metrics:

\begin{itemize}
\tightlist
\item
  For AUC/PR-AUC (involving two independent samples from \(Z=1\) and
  \(Z=0\)):
  \[|AUC_P - AUC_{\tilde{P}}| \leq \rho(\delta_phi) \cdot \left(\frac{1}{\eta} + \frac{1}{1-\eta}\right)\]
\item
  For F1 (joint distribution function, Lipschitz constant \(C_F\)):
  \[|F1_P - F1_{\tilde{P}}| \leq C_F \cdot \rho(\delta_phi)\]
\end{itemize}

\textbf{Step 6}: Combine. \(\square\)

\begin{center}\rule{0.5\linewidth}{0.5pt}\end{center}

\subsection{\texorpdfstring{6 The Role of \(\eta\): How Rare Noise
Amplifies the
Bound}{6 The Role of \textbackslash eta: How Rare Noise Amplifies the Bound}}\label{the-role-of-eta-how-rare-noise-amplifies-the-bound}

The \(\eta\) dependence in the AUC/PR-AUC bounds is a critical
structural feature that merits explicit discussion.

\subsubsection{\texorpdfstring{6.1 Origin of \(\eta\)
Dependence}{6.1 Origin of \textbackslash eta Dependence}}\label{origin-of-eta-dependence}

The amplification factor \(1/\eta + 1/(1-\eta)\) arises from converting
\textbf{marginal} TV bounds to \textbf{conditional} TV bounds:

\[\operatorname{TV}(P(\cdot \mid Z=1), \tilde{P}(\cdot \mid Z=1)) \leq \frac{\operatorname{TV}(P, \tilde{P})}{\mathbb{P}(Z=1)} = \frac{\rho(\delta_phi)}{\eta}\]

This is not an artifact of the proof technique -- it reflects a genuine
statistical phenomenon: when noise is rare (\(\eta \ll 1\)), the
conditional distribution of features given \(Z=1\) is estimated from
very few samples, making it harder to distinguish from the \(Z=0\)
distribution.

\subsubsection{6.2 Amplification Behavior}\label{amplification-behavior}

\begin{longtable}[]{@{}
  >{\centering\arraybackslash}p{(\linewidth - 4\tabcolsep) * \real{0.3333}}
  >{\centering\arraybackslash}p{(\linewidth - 4\tabcolsep) * \real{0.3333}}
  >{\centering\arraybackslash}p{(\linewidth - 4\tabcolsep) * \real{0.3333}}@{}}
\toprule\noalign{}
\begin{minipage}[b]{\linewidth}\centering
Noise rate \(\eta\)
\end{minipage} & \begin{minipage}[b]{\linewidth}\centering
\(1/\eta + 1/(1-\eta)\)
\end{minipage} & \begin{minipage}[b]{\linewidth}\centering
Effective bound multiplier
\end{minipage} \\
\midrule\noalign{}
\endhead
\bottomrule\noalign{}
\endlastfoot
0.50 & 4.00 & 1.00\(\times\) (baseline) \\
0.30 & 4.76 & 1.19\(\times\) \\
0.20 & 6.25 & 1.56\(\times\) \\
0.10 & 11.11 & 2.78\(\times\) \\
0.05 & 21.05 & 5.26\(\times\) \\
0.01 & 101.01 & 25.25\(\times\) \\
0.001 & 1001.00 & 250.25\(\times\) \\
\end{longtable}

As \(\eta \to 0\), the bound diverges:
\(1/\eta + 1/(1-\eta) \to \infty\). This is not a weakness of the bound
but a reflection of the fundamental difficulty: \textbf{rare noise is
information-theoretically harder to detect}.

\subsubsection{6.3 Practical Implications}\label{practical-implications}

\textbf{Diagnostic 1}: If \(\eta < 0.1\) and \(\varepsilon_\phi > 0.3\),
SCX is unlikely to outperform the loss baseline by a meaningful margin.
The AUC bound becomes too loose to guarantee improvement.

\textbf{Diagnostic 2}: When noise is very rare (\(\eta < 0.01\)), do not
rely on AUC/PR-AUC for evaluating SCX. Use F1 instead (which has no
\(\eta\) amplification factor) or re-frame the problem as anomaly
detection.

\textbf{Diagnostic 3}: The F1 bound has \textbf{no} \(\eta\) dependence
because F1 is a function of the joint distribution \(P(\hat{z}, Z)\)
rather than conditional distributions. This makes F1 the preferred
metric in rare-noise regimes.

\textbf{Comparison with Theorem 1}: Theorem 1's F1 bound also involves
\(1/\eta\), but for a different reason (the conversion from per-class
error bounds to F1). In Theorem 1, \(1/\eta\) multiplies the exponential
term and is mitigated by the exponential factor. In Theorem 2, the
\(1/\eta\) factor is additive and not mitigated by any sample size
effect.

\subsubsection{\texorpdfstring{6.4 Boundary Behavior: \(\eta = 0\) and
\(\eta = 1\)}{6.4 Boundary Behavior: \textbackslash eta = 0 and \textbackslash eta = 1}}\label{boundary-behavior-eta-0-and-eta-1}

The theorem statement \textbf{excludes} the boundary cases \(\eta = 0\)
and \(\eta = 1\). This is not a limitation but a reflection of
degeneracy:

\textbf{\(\eta = 0\) (no noise).} If there is no label noise, the noise
detection problem is trivial --- there are no noise samples to detect.
\(P(Z=1) = 0\) means the conditional distribution \(P(\cdot \mid Z=1)\)
is undefined (conditioning on a measure-zero event). The AUC/PR-AUC
bounds, which involve \(1/\eta\), diverge, correctly signaling that
noise detection metrics are not meaningful when noise is absent. In this
case, one should simply use all data for training without noise
filtering.

\textbf{\(\eta = 1\) (all noise).} If all labels are noisy,
\(P(Z=0) = 0\) and the conditional distribution \(P(\cdot \mid Z=0)\) is
undefined. The \(1/(1-\eta)\) term diverges. Again, this is correct: if
all data are noise, there are no clean samples to serve as a reference,
and the noise detection problem is ill-posed. In this case, one should
use robust estimation or outlier detection rather than supervised noise
filtering.

\textbf{Practical note.} In any real application, \(0 < \eta < 1\)
strictly (there is always some noise and always some clean data). The
restriction to the open interval \((0,1)\) is therefore non-restrictive
in practice. The theorem's divergence at the boundaries is not a bug ---
it correctly reflects that the problem becomes degenerate when one class
is absent.

\textbf{Formal adjustment to the proof.} In Step 5 of the proof (§5.3),
the conversion from marginal TV to conditional TV uses:

\[\operatorname{TV}(P(\cdot \mid Z=1), \tilde{P}(\cdot \mid Z=1)) \leq \frac{\operatorname{TV}(P, \tilde{P})}{\eta}\]

This step is valid only when \(\eta > 0\). Similarly, the \(Z=0\)
conditional bound requires \(1-\eta > 0\). The theorem therefore
requires \(\eta \in (0,1)\). At the boundaries, the theorem is vacuous
(the bounds are infinite), which is the correct mathematical behavior
for a degenerate problem.

\begin{center}\rule{0.5\linewidth}{0.5pt}\end{center}

\subsection{\texorpdfstring{7 Estimation of \(\delta_phi\) from
Data}{7 Estimation of \textbackslash delta\_phi from Data}}\label{estimation-of-delta_phi-from-data}

\subsubsection{7.1 k-NN Mutual Information Estimation (Kraskov et al.,
2004)}\label{k-nn-mutual-information-estimation-kraskov-et-al.-2004}

Even without the BBP spectral proxy, the mutual information
\(\delta_phi = I(\phi; S)\) can be estimated directly from data using
the Kraskov-Stoegbauer-Grassberger (KSG) estimator:

\[\hat{I}_{\text{KSG}}(\phi; S) = \psi(K) + \psi(N) - \frac{1}{N} \sum_{i=1}^N \left[\psi(n_{\phi(i)}+1) + \psi(n_{s(i)}+1)\right]\]

where: - \(\psi\) is the digamma function - \(K\) is the number of
nearest neighbors (typically \(K=3\) or \(K=6\)) - \(N\) is the total
number of samples - \(n_{\phi(i)}\) is the number of \(\phi\)-space
points within the same radius as the \(K\)-th neighbor of point \(i\) -
\(n_{s(i)}\) is the number of state-assignment points within the same
radius

\textbf{Practical implementation}: - Requires: feature vectors
\(\{\phi(x_i)\}\) and (proxy) state labels \(\{s_i\}\) - Even
noisy/estimated state labels yield a lower bound on \(\delta_phi\) (by
the data processing inequality) - Python: available via
\texttt{sklearn.feature\_selection.mutual\_info\_classif} (discrete
\(S\)) or \texttt{sklearn.feature\_selection.mutual\_info\_regression}
(continuous \(S\))

\textbf{Caveat}: When state labels are themselves estimated (e.g., via
clustering), \(\hat{I}_{\text{KSG}}(\phi; \hat{S})\) is a \textbf{lower
bound} on \(I(\phi; S)\):

\[I(\phi; \hat{S}) \leq I(\phi; S)\]

by the data processing inequality (\(S \to \phi \to \hat{S}\) is
Markov). Hence \(\hat{I}_{\text{KSG}}(\phi; \hat{S}) \leq \delta_phi\),
and using it in Theorem 2 yields a conservative bound (the true
\(\delta_phi\) may be larger, making the bound looser).

\subsubsection{7.2 Quick Diagnostic via Normalized
Weakness}\label{quick-diagnostic-via-normalized-weakness}

For practitioners who need a quick check without full mutual information
estimation:

\begin{enumerate}
\def\labelenumi{\arabic{enumi}.}
\tightlist
\item
  Cluster \(\phi\) to get estimated states \(\hat{S}\)
\item
  Compute consistency scores \(C(\hat{s})\) for each estimated state
\item
  If \(\operatorname{Var}(\{C(\hat{s})\}) < 0.01\), the feature is
  likely \(\delta_phi\)-weak with \(\varepsilon_\phi > 0.5\)
\item
  For a quantitative estimate, compute \(\hat{I}_{\text{KSG}}\) or use
  the ARI between \(\hat{S}\) and any available metadata
\end{enumerate}

\begin{center}\rule{0.5\linewidth}{0.5pt}\end{center}

\subsection{8 Corollaries}\label{corollaries}

\subsubsection{\texorpdfstring{8.1 Corollary 1: Complete Failure
(\(\delta_phi = 0\))}{8.1 Corollary 1: Complete Failure (\textbackslash delta\_phi = 0)}}\label{corollary-1-complete-failure-delta_phi-0}

If \(\phi(X) \perp S\), then:

\[AUC(h_{\text{SCX}}) = AUC_{\text{base}}, \quad F1(h_{\text{SCX}}) = F1_{\text{base}}\]

If further the noise is loss-uninformative (e.g., uniform):
\(F1_{\text{SCX}} = F1_{\text{rand}}\).

\subsubsection{8.2 Corollary 2: Uniform Label
Noise}\label{corollary-2-uniform-label-noise}

Under uniform label noise, the loss baseline is random:

\[AUC_{\text{base}} = 0.5, \quad PRAUC_{\text{base}} = \eta, \quad F1_{\text{base}} = \frac{2\eta}{1+\eta}\]

Then:

\[\begin{aligned}
AUC(h_{\text{SCX}}) &\leq 0.5 + \rho(\delta_phi) \cdot \left(\frac{1}{\eta} + \frac{1}{1-\eta}\right) \\
F1(h_{\text{SCX}}) &\leq \frac{2\eta}{1+\eta} + C_F \cdot \rho(\delta_phi)
\end{aligned}\]

When \(\delta_phi = 0\), SCX is bounded by the random baseline.

\subsubsection{8.3 Corollary 3: Minimum Information for
Improvement}\label{corollary-3-minimum-information-for-improvement}

To exceed the baseline by \(\Delta\) in F1, the required minimum mutual
information is:

\[\delta_phi_{\min} \geq -\log\!\left(1 - \left(\frac{\Delta}{C_F}\right)^2\right) \approx \left(\frac{\Delta}{C_F}\right)^2 \quad \text{for small } \Delta\]

For AUC, with \(\eta\) dependence:

\[\delta_phi_{\min} \geq -\log\!\left(1 - \left(\Delta_{AUC} \cdot \frac{\eta(1-\eta)}{\min(\eta, 1-\eta)}\right)^2\right)\]

\textbf{Example}: For \(\Delta_{F1} = 0.05\) with \(C_F = 2\):

\[\delta_phi_{\min} \geq -\log(1 - 0.05^2/4) \approx 0.000625 \text{ nats}\]

This is a very small amount of mutual information -- confirming that SCX
can extract value from even weak features, as long as they are not
completely uninformative.

\subsubsection{8.4 Corollary 4: State Count
Interaction}\label{corollary-4-state-count-interaction}

The normalized weakness \(\varepsilon_\phi = \delta_phi / \log K_S\)
determines SCX effectiveness:

\begin{itemize}
\tightlist
\item
  \(\varepsilon_\phi \approx 1\): SCX degrades to baseline
\item
  \(\varepsilon_\phi \approx 0.5\): partial effectiveness
\item
  \(\varepsilon_\phi \approx 0\): strong feature, full SCX capability
\end{itemize}

Larger \(K_S\) demands larger \(\delta_phi\) to maintain the same
\(\varepsilon_\phi\).

\begin{center}\rule{0.5\linewidth}{0.5pt}\end{center}

\subsection{9 Practical Diagnostics}\label{practical-diagnostics}

\subsubsection{9.1 Four Diagnostic Checks}\label{four-diagnostic-checks}

\begin{longtable}[]{@{}
  >{\raggedright\arraybackslash}p{(\linewidth - 6\tabcolsep) * \real{0.2895}}
  >{\raggedright\arraybackslash}p{(\linewidth - 6\tabcolsep) * \real{0.2105}}
  >{\raggedright\arraybackslash}p{(\linewidth - 6\tabcolsep) * \real{0.2895}}
  >{\raggedright\arraybackslash}p{(\linewidth - 6\tabcolsep) * \real{0.2105}}@{}}
\toprule\noalign{}
\begin{minipage}[b]{\linewidth}\raggedright
Diagnostic
\end{minipage} & \begin{minipage}[b]{\linewidth}\raggedright
Method
\end{minipage} & \begin{minipage}[b]{\linewidth}\raggedright
Threshold
\end{minipage} & \begin{minipage}[b]{\linewidth}\raggedright
Action
\end{minipage} \\
\midrule\noalign{}
\endhead
\bottomrule\noalign{}
\endlastfoot
1. Consistency convergence & Variance of \(\{C(\hat{s})\}\) &
\(\operatorname{Var} < 0.01\) & Feature likely too weak \\
2. Random baseline comparison & \(AUC_{\text{base}}\) vs \(0.5\) &
\(AUC_{\text{base}} < 0.55\) & Task itself is hard \\
3. Supervised ARI & ARI(\(\hat{S}\), true \(S\)) if available &
\(< 0.1\) & Feature too weak \\
4. MI estimation & \(\hat{I}_{\text{KSG}}(\phi; \hat{S})\) &
\(\varepsilon_\phi > 0.5\) & Feature too weak \\
\end{longtable}

\subsubsection{9.2 Remedy Directions}\label{remedy-directions}

When \(\varepsilon_\phi > 0.5\):

\begin{enumerate}
\def\labelenumi{\arabic{enumi}.}
\tightlist
\item
  \textbf{Enhance features}: Use stronger descriptors (ACE/SOAP/MACE) or
  error-driven feature learning
\item
  \textbf{Reduce \(K_S\)}: Fewer states reduce the demand on
  \(\delta_phi\)
\item
  \textbf{Abandon state conditioning}: Use global loss threshold
  directly
\item
  \textbf{Reformulate}: Use cross-validation consistency or temporal
  anomaly detection
\end{enumerate}

\begin{center}\rule{0.5\linewidth}{0.5pt}\end{center}

\subsection{10 Relationship to Theorems 1 and
3}\label{relationship-to-theorems-1-and-3}

\begin{longtable}[]{@{}
  >{\raggedright\arraybackslash}p{(\linewidth - 6\tabcolsep) * \real{0.0833}}
  >{\raggedright\arraybackslash}p{(\linewidth - 6\tabcolsep) * \real{0.3056}}
  >{\raggedright\arraybackslash}p{(\linewidth - 6\tabcolsep) * \real{0.3056}}
  >{\raggedright\arraybackslash}p{(\linewidth - 6\tabcolsep) * \real{0.3056}}@{}}
\toprule\noalign{}
\begin{minipage}[b]{\linewidth}\raggedright
\end{minipage} & \begin{minipage}[b]{\linewidth}\raggedright
Theorem 1
\end{minipage} & \begin{minipage}[b]{\linewidth}\raggedright
Theorem 2
\end{minipage} & \begin{minipage}[b]{\linewidth}\raggedright
Theorem 3
\end{minipage} \\
\midrule\noalign{}
\endhead
\bottomrule\noalign{}
\endlastfoot
\textbf{Answer} & When SCX \textbf{can} work & When SCX \textbf{cannot}
work & Why assumptions are \textbf{necessary} \\
\textbf{Key quantity} & \(\Delta_s\), \(\mu_s\), \(M\) &
\(\delta_phi = I(\phi; S)\) & A1-A6 validity \\
\textbf{Condition} & Good state partition & \(\varepsilon_\phi > 0.5\) &
Without A1-A6, unidentifiable \\
\textbf{Guarantee} & Exponential F1 \(\to 1\) & Bounded by baseline
\(+ O(\sqrt{\delta_phi})\) & No guarantee possible \\
\end{longtable}

The three theorems together give:

\[\text{SCX Advantage} \approx I(\phi; S) - O(1/\sqrt{n}) - \text{(assumption violations)}\]

\begin{center}\rule{0.5\linewidth}{0.5pt}\end{center}

\subsection{References}\label{references}

\begin{enumerate}
\def\labelenumi{\arabic{enumi}.}
\tightlist
\item
  Fano, R. M. (1961). \emph{Transmission of Information}. MIT Press.
\item
  Cover, T. M. \& Thomas, J. A. (2006). \emph{Elements of Information
  Theory} (2nd ed.). Wiley.
\item
  Pinsker, M. S. (1964). \emph{Information and Information Stability}.
  Holden-Day.
\item
  Tsybakov, A. B. (2009). \emph{Introduction to Nonparametric
  Estimation}. Springer.
\item
  Kraskov, A., Stoegbauer, H., \& Grassberger, P. (2004). Estimating
  mutual information. \emph{Physical Review E}, 69(6), 066138.
\item
  Bretagnolle, J. \& Huber, C. (1979). Estimation des densites: risque
  minimax. \emph{Zeitschrift fur Wahrscheinlichkeitstheorie}, 47(2),
  119-137.
\end{enumerate}

\begin{center}\rule{0.5\linewidth}{0.5pt}\end{center}

\textbf{Revision notes (2026-06-27)}: 1. \textbf{Bretagnolle-Huber}:
Replaced Pinsker with BH bound; added comparison table and guidance on
selecting the tighter bound. 2. \textbf{\(\delta_phi\) estimation}:
Added Section 7 with k-NN mutual information estimation (Kraskov et al.,
2004). 3. \textbf{\(\eta\) dependence}: Added Section 6 analyzing how
rare noise amplifies the bound; included practical diagnostics. 4.
\textbf{Presentation}: Unified notation; added explicit cross-theorem
dependency mapping.

\end{document}
